\section*{Chapter \ref{ch:m2:reading-a-macro-calendar-and-a-rates-screen} --- Reading a Macro Calendar and a Rates Screen}

\begin{solution}{exo:m2:reading-a-macro-calendar-and-a-rates-screen:1}
From midnight on Saturday 17 October to 23:59 on Thursday 29 October 2026.
\end{solution}

\begin{solution}{exo:m2:reading-a-macro-calendar-and-a-rates-screen:2}
\omterm{def:m2:reading-a-macro-calendar-and-a-rates-screen:bull}{Bull steepening}: yields fell, the 2-year more, so 2s10s rose by 6 basis points.
\end{solution}

\begin{solution}{exo:m2:reading-a-macro-calendar-and-a-rates-screen:3}
2s10s $415 - 361 = 54$, 5s30s $477 - 371 = 106$, fly $2 \times 371 - 361 - 415 = -34$ basis
points.
\end{solution}

\begin{solution}{exo:m2:reading-a-macro-calendar-and-a-rates-screen:4}
The 2-year yield is close to the average policy rate expected over two years, and the
report changes that expectation most directly; the 10-year also contains expected rates
further out, which move less with one month's data, and a term premium.
\end{solution}

\begin{solution}{exo:m2:reading-a-macro-calendar-and-a-rates-screen:5}
$50\,000 \times (-5 + 15) = \text{USD}~500\,000$; a parallel fall of 10 earns nothing to
first order.
\end{solution}

\begin{solution}{exo:m2:reading-a-macro-calendar-and-a-rates-screen:6}
Because the outcome is close to a coin toss for anyone without an edge on the number, the
moves are two to three times larger than usual, and liquidity thins around the release:
the risk rises more than the expected reward.
\end{solution}

\begin{solution}{exo:m2:reading-a-macro-calendar-and-a-rates-screen:7}
The standard error is about the noise over the square root of the number of releases,
$3/\sqrt n$ with unit-variance surprises: 0.5 needs about 36 releases, three years of monthly
data. Simulated with 36 releases, the average standard error is 0.51.
\end{solution}

\begin{solution}{exo:m2:reading-a-macro-calendar-and-a-rates-screen:8}
The 10 basis points are the average absolute move of the 2-year, not a bad day, and they
say nothing about the position's other exposures: a curve position's risk is the slope's
move times its DV01, a portfolio's the combination of its exposures along the curve. Use
the distribution of event-day moves of each exposure, including its tail: the 2-year moved
up to 30 basis points on payroll days in the sample.
\end{solution}

\begin{solution}{pb:m2:reading-a-macro-calendar-and-a-rates-screen:1}
\textbf{1.} 0.0183 for the 2-year and 0.0772 for the 10-year, per 100.
\textbf{2.} USD~77\,168 per basis point.
\textbf{3.} USD~421 million of the 2-year.
\textbf{4.} Because the 10-year moves the price four times as much per basis point:
equal notional would be mostly a bet on the level of the 10-year.
\textbf{5.} Nothing, to first order.
\textbf{6.} USD~694\,512 to first order; USD~688\,510 fully repriced.
\textbf{7.} A loss of USD~771\,680 to first order; USD~772\,969 repriced.
\textbf{8.} \omterm{def:m2:reading-a-macro-calendar-and-a-rates-screen:bull}{Bull steepening}, then \omterm{def:m2:reading-a-macro-calendar-and-a-rates-screen:bull}{bear flattening}.
\textbf{9.} 34 times out of 44.
\textbf{10.} Convexity and the different sizes of the moves: the 2-year and 10-year prices
are not linear in their yields, and the 10-year's convexity is larger.
\textbf{11.} No: it begins the next day, Saturday 5 September.
\textbf{12.} Policymakers cannot comment on it before the meeting, so the market's reading
of the CPI is not corrected or confirmed by them until the decision.
\textbf{13.} 10.2 basis points against 4.6, 2.2 times as much.
\textbf{14.} The CPI on 11 September, Treasury auctions, and anything else scheduled, such
as speeches before the blackout begins.
\textbf{15.} The largest payroll-day move in the 2s10s slope in the sample was 11 basis
points (1 August 2025): a DV01 of about USD~45\,000 per leg keeps that loss within
USD~500\,000; the trader's USD~77\,000 would lose about USD~850\,000 on such a day against it.
\textbf{16.} Because a weak report may move the long end as much, if it raises doubts about
term premium or supply, or because the market had already priced it; and because the
report is only one of the things that moved yields that day.
\textbf{17.} The 2-year, the \omterm{def:m2:short-term-interest-rate-futures:path}{implied policy path} and 2s10s before and after; the reaction
of the long end; volumes and liquidity; the next events on the calendar.
\textbf{18.} The \omterm{def:m2:reading-a-macro-calendar-and-a-rates-screen:steep}{steepener} profits when the market cuts its expected policy path at the
front: it is a way to bet on the path of chapter 8 without holding the level of rates.
\textbf{19.} Named result: \emph{payrolls Friday}: buy USD~421 million of the 2-year
against USD~100 million of the 10-year, USD~77\,168 per basis point on each leg; the
\omterm{def:m2:reading-a-macro-calendar-and-a-rates-screen:steep}{steepener} gains USD~694\,512 on a repeat of 2 August 2024 (688\,510 repriced) and loses
USD~771\,680 on a repeat of 4 October 2024 (772\,969 repriced).
\textbf{20.} On the slope between the two maturities, not on the level of rates.
\end{solution}

\begin{solution}{iq:m2:reading-a-macro-calendar-and-a-rates-screen:1}
The employment report and the CPI, at 08:30 Eastern time, the first usually on the first
Friday of the month and the second in the middle of it; the FOMC decisions after its eight
scheduled meetings; also retail sales, GDP, PCE inflation and Treasury refunding and
auctions.
\iqlookfor{the main releases and their timing.}
\end{solution}

\begin{solution}{iq:m2:reading-a-macro-calendar-and-a-rates-screen:2}
\omterm{def:m2:reading-a-macro-calendar-and-a-rates-screen:bull}{Bull steepening}: yields fall, the front end more, as when weak data pull expected policy
rates down. \omterm{def:m2:reading-a-macro-calendar-and-a-rates-screen:bull}{Bear flattening}: yields rise, the front end more, as when strong data or a
hawkish central bank push expected policy rates up.
\iqlookfor{direction, shape, and the policy-path channel.}
\end{solution}

\begin{solution}{iq:m2:reading-a-macro-calendar-and-a-rates-screen:3}
Buy the 2-year and sell the 10-year in DV01-equal amounts, in bonds, futures or swaps. It
gains when 2s10s widens. Risks: the slope moving the other way, non-parallel moves within
each leg, financing and roll (the carry of each leg), and basis between the instruments.
\iqlookfor{DV01 neutrality and the remaining risks.}
\end{solution}

\begin{solution}{iq:m2:reading-a-macro-calendar-and-a-rates-screen:4}
Collect the release, the consensus and the market's move in a narrow window around the
release time; standardise the surprise; regress the moves on the surprises across many
releases; check the stability of the coefficient over time and whether large surprises
act linearly.
\iqlookfor{surprise, narrow window, regression, stability.}
\end{solution}

\begin{solution}{iq:m2:reading-a-macro-calendar-and-a-rates-screen:5}
The curve levels and changes at key maturities, slopes and a fly; the \omterm{def:m2:short-term-interest-rate-futures:path}{implied policy path}
and its change; \omterm{def:m2:interest-rate-swaps:spread}{swap spreads} and the \omterm{def:m2:fx-swaps-forwards-and-the-cross-currency-basis:basis}{cross-currency basis}; implied volatility; financing
rates; the day's calendar and the positions' DV01 by bucket against their limits.
\iqlookfor{level, shape, expectations, relative value, risk, calendar.}
\end{solution}

\begin{solution}{iq:m2:reading-a-macro-calendar-and-a-rates-screen:6}
Subscribe to the release feed with the lowest latency available and parse the numbers
into standard fields; compute surprises against pre-loaded consensus; push the new
figures to the pricing engines, which recompute curves and quotes incrementally; limit
checks and risk refresh in the same event loop; everything time-stamped for replay; kill
switches if the feed is malformed or late.
\iqlookfor{pre-loaded context, incremental recompute, fail-safe.}
\end{solution}
